\chapter{Diverticulitis}

\section*{Definition and Overview}
Diverticulitis is inflammation or infection of diverticula, the small pouches that can form in the wall of the large intestine. It is a complication of diverticular disease, in which a pouch becomes blocked or irritated, allowing bacteria to multiply and cause inflammation. Diverticulitis typically causes constant pain, usually in the lower left abdomen, along with fever, nausea, and changes in bowel habit. Most cases are mild and can be treated at home with rest, a liquid diet, and sometimes antibiotics, but severe cases require hospital care and, rarely, surgery. This chapter explains the causes, symptoms, diagnosis, and treatment of diverticulitis.

\section*{Epidemiology}
Diverticulitis develops in around 10--25\% of people with diverticulosis, and it becomes more common with age, affecting mainly people over 50. It is one of the most common reasons for hospital admission with a digestive condition in Australia. The rate of diverticulitis has risen in recent decades, possibly related to low-fibre diets and ageing. Most attacks are uncomplicated and resolve with conservative treatment, with around 15--30\% of people experiencing a recurrence.

\section*{Aetiology and Pathophysiology}
\begin{itemize}
    \item \textbf{Diverticula:} pouches form where the bowel wall is weak, driven by pressure from a low-fibre diet and constipation
    \item \textbf{Inflammation:} stool or food becomes trapped in a pouch, and bacteria overgrow
    \item \textbf{Microperforation:} inflammation can cause small perforations, with infection spreading into surrounding tissue
    \item \textbf{Complications:} abscess, larger perforation, fistula, obstruction, and bleeding
    \item \textbf{Risk factors:} age, low-fibre diet, obesity, smoking, and certain medicines such as NSAIDs
    \item \textbf{Mechanism:} the trapped material inflames the pouch, causing localised pain and systemic symptoms
\end{itemize}

\section*{Clinical Features}
\begin{itemize}
    \item Constant abdominal pain, most often in the lower left side
    \item Tenderness over the affected area
    \item Fever, chills, and a general feeling of being unwell
    \item Nausea, vomiting, and loss of appetite
    \item A change in bowel habit, often constipation
    \item Urgency and frequent urination when the bladder is irritated
    \item Severe cases: a rigid abdomen, severe tenderness, and signs of peritonitis
\end{itemize}

\section*{Differential Diagnosis}
\begin{itemize}
    \item Irritable bowel syndrome and functional abdominal pain
    \item Inflammatory bowel disease, including Crohn's disease and ulcerative colitis
    \item Bowel cancer, which can present similarly
    \item Appendicitis, when inflammation is on the right side
    \item Gynaecological conditions, including ovarian cysts and pelvic inflammatory disease
    \item Kidney stones and urinary tract infection
    \item Ischaemic colitis, from reduced blood flow to the bowel
\end{itemize}

\section*{When to Seek Medical Care}
See a doctor for persistent abdominal pain, fever, or changes in bowel habit. Anyone with suspected diverticulitis needs assessment to confirm the diagnosis and decide whether home or hospital care is needed.

\textbf{Call 000 (or local emergency services) for:}
\begin{itemize}
    \item Sudden, severe abdominal pain
    \item High fever with abdominal pain and vomiting
    \item A rigid, tender abdomen with signs of peritonitis
    \item Significant bleeding from the bowel
    \item Inability to keep down fluids or signs of dehydration
\end{itemize}

\section*{Diagnosis and Investigations}
\begin{itemize}
    \item \textbf{History and examination:} the pattern of pain, fever, and tenderness
    \item \textbf{Blood tests:} white cell count and inflammatory markers
    \item \textbf{CT scan:} the most accurate test for diagnosing diverticulitis and complications
    \item \textbf{Ultrasound:} an alternative in some settings
    \item \textbf{Urine tests:} to exclude urinary causes
    \item \textbf{Colonoscopy:} performed after recovery to exclude other conditions and assess the bowel
\end{itemize}

\section*{Management}
\begin{itemize}
    \item \textbf{Mild diverticulitis at home:} rest, fluids, a clear or low-residue diet, and pain relief
    \item \textbf{Antibiotics:} used for moderate or complicated disease; not always needed for mild cases
    \item \textbf{Hospital care:} intravenous fluids and antibiotics for severe or complicated disease
    \item \textbf{Abscess drainage:} image-guided drainage of collections
    \item \textbf{Surgery:} for perforation, recurrent severe attacks, or complications such as fistulas
    \item \textbf{Pain relief:} paracetamol is preferred; anti-inflammatory medicines are generally avoided during an attack
    \item \textbf{Follow-up:} review, a gradual return to a high-fibre diet, and colonoscopy after recovery
\end{itemize}

\section*{Complications}
\begin{itemize}
    \item Abscess formation
    \item Perforation with peritonitis, a surgical emergency
    \item Fistulas to the bladder, vagina, or skin
    \item Bowel obstruction from scarring
    \item Significant bleeding
    \item Recurrent attacks
\end{itemize}

\section*{Prognosis and Outcomes}
Most people with a first attack of diverticulitis recover fully with simple treatment, and most do not need surgery. Around 15--30\% have recurrent attacks, but the risk falls with a high-fibre diet, healthy weight, and smoking cessation. Serious complications affect a minority, and emergency surgery is needed in only a small proportion. The long-term outlook for diverticulitis is good with appropriate management.

\section*{Prevention and Self-Care}
\begin{itemize}
    \item Eat a high-fibre diet to prevent diverticula and reduce recurrence
    \item Drink plenty of water
    \item Exercise and maintain a healthy weight
    \item Do not smoke
    \item Limit anti-inflammatory medicines, which can worsen diverticulitis
    \item Follow the doctor's advice on diet during and after an attack
    \item See a doctor promptly for abdominal pain with fever
\end{itemize}

\section*{Sources and Further Reading}
The following reputable sources provide further information for healthcare students and patients seeking reliable, current advice about diverticulitis.

\begin{itemize}
    \item The Gut Foundation. \textit{Diverticulitis: information and management.} https://www.gutfoundation.com.au/
    \item Healthdirect Australia. \textit{Diverticulitis: symptoms and treatment.} https://www.healthdirect.gov.au/
    \item NHS. \textit{Diverticulitis: treatment and recovery.} https://www.nhs.uk/conditions/diverticular-disease-and-diverticulitis/
    \item Mayo Clinic. \textit{Diverticulitis: diagnosis and treatment.} https://www.mayoclinic.org/
    \item MSD Manual. \textit{Diverticulitis: professional overview.} https://www.msdmanuals.com/
    \item MedlinePlus. \textit{Diverticulitis.} https://medlineplus.gov/
\end{itemize}
